% ===== PRÉAMBULE CANONIQUE — à coller VERBATIM en tête de CHAQUE .tex =====
\documentclass[11pt,a4paper]{article}
\usepackage[utf8]{inputenc}\usepackage[T1]{fontenc}\usepackage{lmodern}
\usepackage[french]{babel}
\usepackage{amsmath,amssymb,mathtools}\usepackage{icomma}
\usepackage{siunitx}\sisetup{locale=FR}  % \qty{}{}, \si{}, \ang{}
\usepackage{xcolor}\usepackage{colortbl}
\usepackage{tikz}\usepackage{pgfplots}\pgfplotsset{compat=1.18}
\usepackage[siunitx,europeanresistors]{circuitikz} % circuits (élec.)
\usepackage[most]{tcolorbox}\usepackage{enumitem}\usepackage{array,booktabs}
\usepackage[a4paper,margin=2.2cm]{geometry}
\usetikzlibrary{arrows.meta,calc,angles,quotes,positioning,patterns,%
  backgrounds,shapes.geometric,decorations.pathmorphing,decorations.markings,intersections}
\definecolor{primary}{RGB}{222,49,99}\definecolor{primarydark}{RGB}{150,22,55}
\definecolor{defcol}{RGB}{0,114,178}\definecolor{methcol}{RGB}{52,92,148}
\definecolor{excol}{RGB}{0,135,90}\definecolor{attncol}{RGB}{200,40,55}
\tcbset{boxbase/.style={enhanced,breakable,boxrule=0.9pt,arc=2.5pt,
  left=3mm,right=3mm,top=2.3mm,bottom=2.3mm,fonttitle=\bfseries,coltitle=white}}
% --- boîtes TUTORIEL (noms EXACTS, ne pas renommer) ---
\newtcolorbox{cle}[1][]{boxbase,colback=defcol!6,colframe=defcol,
  title={\textsc{À retenir}\ifx\relax#1\relax\relax\else\textmd{ : \itshape #1}\fi}}
\newtcolorbox{methode}[1][]{boxbase,colback=methcol!7,colframe=methcol,sharp corners,
  title={\textsc{Méthode}\ifx\relax#1\relax\relax\else\textmd{ --- \itshape #1}\fi}}
\newtcolorbox{exemplebox}[1][]{boxbase,colback=excol!5,colframe=excol,
  title={\textsc{Exemple}\ifx\relax#1\relax\relax\else\textmd{ : \itshape #1}\fi}}
\newtcolorbox{piege}[1][]{boxbase,colback=attncol!6,colframe=attncol,
  title={\textsc{Piège}\ifx\relax#1\relax\relax\else\textmd{ : \itshape #1}\fi}}
\newtcolorbox{remarque}[1][]{boxbase,colback=black!4,colframe=black!45,
  title={\textsc{Remarque}\ifx\relax#1\relax\relax\else\textmd{ : \itshape #1}\fi}}
% --- boîtes EXERCICE / CORRIGÉ (noms EXACTS, auto-comptées) ---
\newtcolorbox[auto counter]{exo}[1][]{boxbase,colback=primary!4,colframe=primary,
  title={\textsc{Exercice}~\thetcbcounter\ifx\relax#1\relax\relax\else\textmd{ --- \itshape #1}\fi}}
\newtcolorbox[auto counter]{corrige}[1][]{boxbase,colback=excol!4,colframe=excol,
  title={\textsc{Corrigé}~\thetcbcounter\ifx\relax#1\relax\relax\else\textmd{ --- \itshape #1}\fi}}
\newcommand{\vv}[1]{\vec{#1}}\newcommand{\ds}{\displaystyle}
\newcommand{\R}{\mathbb{R}}\newcommand{\N}{\mathbb{N}}\newcommand{\Z}{\mathbb{Z}}
% (un \usepackage{fancyhdr}+en-tête est possible ; il est ignoré côté web)
% =========================================================================
\begin{document}

\begin{center}
{\large\bfseries\color{primarydark} Thème 4 — Niveau Facile}\\[-2pt]
{\color{primary}\rule{5cm}{1pt}}
\end{center}

\begin{exo}[reconnaître la diffraction]
Des vagues planes (longueur d'onde $\lambda$) arrivent sur une barrière percée d'une petite
ouverture de taille comparable à $\lambda$. Après l'ouverture, on observe la figure suivante.
\begin{center}
\begin{tikzpicture}[scale=0.95]
  \foreach \x in {0,0.55,1.1,1.65,2.2}{\draw[defcol,line width=1pt](\x,-1.8)--(\x,1.8);}
  \draw[->,>=Stealth,black](0.9,-2.15)--(2.0,-2.15)node[right,font=\footnotesize]{$\vv v$};
  \draw[black,line width=2.2pt](3.2,0.4)--(3.2,2.0);
  \draw[black,line width=2.2pt](3.2,-2.0)--(3.2,-0.4);
  \foreach \r in {0.5,1.0,1.5,2.0,2.5}{
     \draw[defcol,dashed,line width=0.8pt]([shift={(3.2,0)}]72:\r) arc (72:-72:\r);}
\end{tikzpicture}
\end{center}
\begin{enumerate}[label=\alph*)]
  \item Comment s'appelle ce phénomène (l'étalement de l'onde après l'ouverture) ?
  \item Serait-il aussi net si l'ouverture était \emph{très grande} devant $\lambda$ ? Justifie.
\end{enumerate}
\end{exo}

\begin{exo}[constructive ou destructive ?]
Deux ondes cohérentes de longueur d'onde $\lambda=\SI{2}{\metre}$ se superposent en un point. On
donne la différence de marche $\delta$.
\begin{enumerate}[label=\alph*)]
  \item $\delta=\SI{6}{\metre}$ : l'interférence est-elle constructive ou destructive ?
  \item $\delta=\SI{5}{\metre}$ : même question.
\end{enumerate}
\end{exo}

\begin{exo}[de la différence de marche au déphasage]
On rappelle $\Delta\varphi=2\pi\,\delta/\lambda$.
\begin{enumerate}[label=\alph*)]
  \item Pour $\delta=\lambda/2$, calcule $\Delta\varphi$. Quelle interférence obtient-on ?
  \item Pour $\delta=\lambda$, calcule $\Delta\varphi$. Quelle interférence obtient-on ?
\end{enumerate}
\end{exo}

\begin{exo}[calculer l'interfrange]
Un dispositif de fentes de Young utilise une lumière de longueur d'onde
$\lambda=\SI{500}{\nano\metre}$, deux fentes distantes de $a=\SI{0.2}{\milli\metre}$, un écran à
$D=\SI{1}{\metre}$.
\begin{enumerate}[label=\alph*)]
  \item Convertis $\lambda$ et $a$ en mètres.
  \item Calcule l'interfrange $i=\lambda D/a$.
\end{enumerate}
\end{exo}

\begin{exo}[position des franges brillantes]
Sur un écran, l'interfrange mesuré vaut $i=\SI{2.5}{\milli\metre}$. Les franges brillantes sont
en $x_n=n\,i$ ($n=1,2,3,\dots$), comptées depuis la frange centrale.
\begin{enumerate}[label=\alph*)]
  \item Donne la position $x_1$ de la 1\textsuperscript{re} frange brillante.
  \item Donne $x_2$.
  \item Donne $x_3$.
\end{enumerate}
\end{exo}

\end{document}
